\chapter{ACRS with Open-Weight Agents on AppWorld}
\label{ch:appworld}

\begin{chaptersummary}
This chapter tests ACRS where it matters most: with real open-weight language models acting inside
applications whose state they change. Each objective is specialised to AppWorld's nine applications and
457 APIs, the system is run on a single free or low-cost GPU, and it is scored by the benchmark's own
unit tests, by collateral-damage tests, and by paired significance tests against a reference agent on
the same model.
\end{chaptersummary}

\section{Why AppWorld}
\label{sec:aw-why}

The ToolBench and KCC layers measure the orchestration mechanisms with simulated agents. AppWorld
\citep{trivedi2024appworld} removes that limitation. Its tasks require long sequences of API calls across
applications such as a payment app, a music app, a file system and e-mail; agents write Python that
calls the APIs; and each task is graded by unit tests on the final state. Some of those tests pass when
the agent does nothing and fail when it changes something it was not asked to change, which turns
collateral damage into a measured quantity. The benchmark therefore tests all four objectives at once:
API retrieval from a large catalogue inside a long episode (O1), the choice of model per step (O2),
whether a state-changing action can be trusted before it executes (O3), and how much of a 30--60-step
history to keep (O4).

\section{Specialisation of the Objectives}
\label{sec:aw-spec}

\paragraph{Side effects.} AppWorld's API documentation records each endpoint's HTTP method. A call is
read-only if its method is GET and state-changing otherwise; this classification is exact, so O3 spends
its effort only on steps that change the world, and the policy stage can single out deletions.

\paragraph{O1.} Before each step the retriever ranks API documentation by BM25 over the instruction and
the last two code blocks, boosted by a decayed memory of the applications the episode has used. When an
application enters the episode its login API is added automatically, because protected calls need that
application's access token. APIs the episode has committed to stay visible. Only the selected
documentation (ten APIs by default) enters the prompt, so a small model works within a fixed context
budget. The analogue retrieves once from the instruction and never updates.

\paragraph{O2.} A two-arm contextual Thompson sampler chooses the small or the large model for each step.
The context records whether the last step failed, whether the step is likely to change state and whether
the episode is long. The small model is penalised on risky contexts until it has proved itself there,
which is the risk-signed term, and the large model carries a cost penalty. The reward is an error-free,
policy-allowed step. The analogue always uses the large model.

\paragraph{O3.} When a proposed code block contains a state-changing call, ACRS samples further
candidate programs, saves the world state, executes every candidate from that state, records the
state-changing calls each made and whether it raised an error, and restores the state after each. This is
the strategy that strong ARC solvers use for abstract reasoning: generate several programs, run them, and
keep the result that independent programs agree on \citep{chollet2019measure}. A candidate's vote
weight is one if it executed cleanly; the outcome signature with the largest weighted share is
certified if the share reaches $\theta$ with at least two supporters, and only then is it executed for
real. Otherwise MNCD abstains: nothing is executed, the agent is told that the candidates disagreed, and
the router escalates. The analogue executes the single greedy proposal.

\paragraph{G1.} Before anything executes, a policy check denies deletions the instruction does not ask
for and writes to applications outside the task's scope, and tells the agent why. It is the AppWorld
form of approve-before-execute.

\paragraph{O4.} The history governor keeps the task, the last three steps and every output line that
carries an access token, an identifier or an error, and compresses the rest head-and-tail. A PI loop
adjusts the compression length so that the share of such critical lines retained stays above 0.9 while
the history stays under its token budget. The analogue keeps only the most recent steps.

\section{Experimental Setup}
\label{sec:aw-setup}

\begin{itemize}[leftmargin=*]
\item \textbf{Models.} Qwen3 open-weight models \citep{yang2025qwen3} served with vLLM
\citep{kwon2023vllm}, chosen by GPU memory: a 32B (A100), 14B (L4) or 8B (T4) model in 4-bit AWQ as the
large tier, and a 1.7B or 4B model as the small tier. Thinking mode is disabled.
\item \textbf{Configurations.} A ReAct-style reference agent on the same models, full ACRS, and
leave-one-out ablations without O1, O2, O3, O4 or G1.
\item \textbf{Protocol.} Development and ablations on \texttt{dev}; final configurations once on
\texttt{test\_normal} and \texttt{test\_challenge} with three seeds; no inspection of test tasks.
\item \textbf{Budget.} At most 40 steps per task and a token ceiling per task; candidate executions in O3
count toward AppWorld's interaction limit and are reported.
\item \textbf{Metrics.} Official TGC and SGC, collateral-damage rate, tokens, router mix, consensus
certifications and abstentions, policy denials, and mean history fidelity; McNemar's exact test and a
paired bootstrap interval on per-task success.
\end{itemize}

\section{Results}
\label{sec:aw-results}

% AppWorld results are inserted from experiments/results/appworld/*.summary.json once the GPU runs finish.
% No AppWorld number is reported anywhere in this thesis before then.
\textit{The development-split pilot and ablations are running on the free and Colab GPUs described above;
this section reports their outcomes, and the test-split results, from the run summaries once they are
complete.}

\section{Position Relative to the Leaderboard}
\label{sec:aw-leaderboard}

The strongest AppWorld entries (November 2026) complete 91--94\,\% of \texttt{test\_normal} tasks with
frontier models, and the best entry using a model of at most 14B parameters reaches 86.9\,\%. ACRS on a
free or low-cost GPU is compared primarily with the reference agent on the \emph{same} model, which
isolates the contribution of the architecture, and secondarily with leaderboard entries of similar model
size. Frontier entries are cited as an upper reference, not as the claim. ACRS also reports cost,
collateral damage and certified error, which no current entry reports.
